% Capítulo — TCC 1
\chapter{Introdução}
\label{chap:introducao}

% Abertura do capítulo: um parágrafo dizendo o que ele apresenta.

\section{Contexto}
\label{sec:introducao-contexto}
% Dados públicos federais em silos, com baixa interoperabilidade semântica; o GovHub-br e a parceria UnB--ministérios.

\section{Problema}
\label{sec:introducao-problema}
% Sem chave comum declarada, cruzar duas bases exige descobrir a chave à mão (códigos, nomes e formatos diferentes), com conhecimento implícito do analista.

\section{Justificativa}
\label{sec:introducao-justificativa}
% Custo de descobrir a chave a cada novo par de bases; o padrão do SPAPI-Tester (LLM embutido no processo) ainda não aplicado a dados abertos; SDD com agentes de IA ainda pouco avaliado empiricamente.

\section{Questões de pesquisa}
\label{sec:introducao-questoes-de-pesquisa}
% QP1 a QP5; tabela com pergunta e como será respondida.

\section{Objetivos}
\label{sec:introducao-objetivos}

\subsection{Objetivo geral}
\label{subsec:introducao-objetivo-geral}
% Propor, implementar por Spec-Driven Development e avaliar um pipeline que use modelos de linguagem embutidos entre estágios determinísticos para descobrir chaves de integração entre bases governamentais heterogêneas.

\subsection{Objetivos específicos}
\label{subsec:introducao-objetivos-especificos}
% Lista numerada: Decision Layer com DSPy (Abstenção, Catálogo de Transformações), estágio Jinja → dbt, benchmark (Pares Reais e Perturbados), avaliação contra baselines (QP1–QP4), construção por SDD com comparação pareada (QP5).

\section{Contribuições esperadas}
\label{sec:introducao-contribuicoes-esperadas}
% Pipeline aberto de descoberta de chaves com LLM embutido, benchmark de chaves por categoria de atrito, replicação do artigo-base, evidência pareada do efeito do SDD.

\section{Organização do trabalho}
\label{sec:introducao-organizacao-do-trabalho}
% Um parágrafo por capítulo.
